\section{МОДЕЛЬ ТРАНСФОРМАЦІЇ НАВЧАЛЬНИХ ПРОГРАМ ДЛЯ ФОРМУВАННЯ ІКК}\label{sec:3}\label{chap3}

У попередніх розділах дисертації здійснено теоретичний аналіз проблеми формування інформаційно-комунікативної компетентності майбутніх дизайнерів (розділ~\ref{sec:1}) та комплексний аналіз стану формування ІКК в українській дизайн-освіті на основі даних (розділ~\ref{sec:2}). Результати аналізу виявили: (а)~системну прогалину~--- 0\% покриття ШІ-тематики у 122 акредитованих програмах спеціальності 022~Дизайн; (б)~нисхідний тренд ІКК-змісту в ОПП КДПУ ($-$35\% за 2018--2021~рр.); (в)~незадоволеність випускників рівнем ІКТ-підготовки (до 33,3\%); (г)~пряму рекомендацію роботодавців щодо інтеграції ШІ; (д)~суттєвий розрив із міжнародними програмами, що мають від 1 до 4 ШІ-дисциплін.

Третій розділ присвячено розробці моделі трансформації навчальних програм для формування ІКК майбутніх дизайнерів. Спираючись на інтегративну теоретичну рамку (п.~\ref{subsec:1.4}) та емпіричні дані (розділ~\ref{sec:2}), обґрунтовано структурно-функціональну модель, розроблено рекомендації щодо трансформації навчальних дисциплін та визначено педагогічні умови впровадження.

\subsection{Структурно-функціональна модель формування ІКК}\label{subsec:3.1}

Розробка моделі формування інформаційно-комунікативної компетентності майбутніх дизайнерів ґрунтується на результатах теоретичного аналізу (розділ~\ref{sec:1}) та емпіричних даних (розділ~\ref{sec:2}), що обґрунтовують необхідність системної трансформації навчальних програм спеціальності 022~Дизайн. Модель побудовано як структурно-функціональну систему, що відображає цілі, теоретико-методологічне підґрунтя, зміст, організаційно-технологічне забезпечення та діагностико-результативний механізм формування ІКК.

Методологічною основою моделювання є системний підхід, що передбачає розгляд освітнього процесу як цілісної системи взаємопов'язаних компонентів. Структурно-функціональне моделювання є усталеним методом педагогічних досліджень, що дозволяє представити складний педагогічний процес у вигляді ієрархічно організованої системи блоків з визначеними функціями та зв'язками між ними. У контексті нашого дослідження модель виконує подвійну функцію: теоретичну (систематизація та інтеграція результатів аналізу) та практичну (керівництво для трансформації навчальних програм).

Структурно-функціональна модель формування ІКК майбутніх дизайнерів складається з п'яти взаємопов'язаних блоків: мотиваційно-цільового, теоретико-методологічного, змістового, організаційно-технологічного та діагностично-результативного (рис.~\ref{fig:structural-model}).

\subsubsection{Мотиваційно-цільовий блок}\label{sssec:3.1.1}

Мотиваційно-цільовий блок визначає соціальне замовлення на підготовку дизайнерів з розвиненою ІКК та формулює мету моделі. Цей блок є вихідним, оскільки саме зовнішні та внутрішні чинники обумовлюють необхідність трансформації навчальних програм.

\textbf{Соціальне замовлення} формується на перетині трьох потоків:

\begin{enumerate}
    \item \textit{Ринкова потреба}. Опитування роботодавців випускників КДПУ (п.~\ref{subsec:2.2}) виявило, що роботодавці безпосередньо рекомендують <<використання штучного інтелекту у сфері дизайну>>. Оцінка ІКТ-навичок випускників коливається від 3,5 до 4,0 балів із 5~--- нижче середнього порівняно з іншими якостями. Ці дані підтверджують наявність розриву між підготовкою та очікуваннями ринку праці.

    \item \textit{Незадоволеність здобувачів}. Опитування 32 випускників трьох програм КДПУ засвідчило: 33,3\% випускників програми <<Графічний дизайн>> не задоволені рівнем ІКТ-підготовки~--- це найвищий показник незадоволеності серед усіх 9 аспектів якості. У програмі <<Дизайн середовища>> 85,7\% випускників оцінили рівень ІКТ на 1--3 бали з 5. Студентське опитування 2021~р. ($n = 61$) виявило статистично значущу незадоволеність матеріально-технічним забезпеченням ($H = 8{,}708$, $p = 0{,}013$), особливо у програмі ГД, яка є найбільш цифрово-залежною.

    \item \textit{Нормативні вимоги}. Стандарт В2 Дизайн визначає компетентності ФК8 (здатність застосовувати комп'ютерні технології) та ФК12 (здатність використовувати сучасне ПЗ), проте не згадує ШІ-інструменти, промпт-інженерію чи генеративний дизайн (таблиця~\ref{tab:b2-ikk-mapping}). Європейська рамка DigComp~2.2 вже визнала ШІ-грамотність складовою цифрової компетентності, що створює нормативний тиск на оновлення стандартів.
\end{enumerate}

\textbf{Мета моделі}: забезпечення формування ІКК майбутніх дизайнерів шляхом системної трансформації навчальних програм з інтеграцією ШІ-інструментів, оновленням змісту за 4 компонентами ІКК та впровадженням адаптованих методів студійної педагогіки.

\textbf{Завдання моделі} конкретизують мету за 4 компонентами ІКК:
\begin{itemize}
    \item формування \textbf{інформаційно-аналітичного} компоненту: здатність до пошуку, критичної оцінки та синтезу професійної інформації, у т.\,ч. згенерованої ШІ;
    \item формування \textbf{комунікативного} компоненту: здатність до ефективної візуальної та вербальної комунікації за допомогою цифрових засобів і ШІ-інструментів;
    \item формування \textbf{технологічного} компоненту: здатність застосовувати сучасні цифрові та ШІ-інструменти дизайну (промпт-інженерія, генеративний дизайн, ШІ-робочі процеси);
    \item формування \textbf{рефлексивного} компоненту: здатність до самоаналізу, критичного оцінювання результатів ШІ-генерації та етичної рефлексії щодо використання технологій.
\end{itemize}

\subsubsection{Теоретико-методологічний блок}\label{sssec:3.1.2}

Теоретико-методологічний блок визначає науково-теоретичне підґрунтя моделі та дидактичні принципи, що забезпечують її реалізацію. Цей блок спирається на інтегративну теоретичну рамку, обґрунтовану в п.~\ref{subsec:1.4}.

\textbf{Теоретичні підходи.} Модель ґрунтується на трьох взаємопов'язаних теоретичних перспективах:

\begin{enumerate}
    \item \textit{Модель TPACK}, адаптована для дизайн-освіти~\cite{Mishra2006}. У контексті формування ІКК дизайнерів компоненти TPACK набувають специфічного змісту: технологічні знання (TK) включають як традиційні цифрові інструменти, так і ШІ-інструменти нового покоління; педагогічні знання (PK) охоплюють студійні методи, адаптовані для роботи з ШІ; змістові знання (CK) інтегрують класичні дизайнерські дисципліни з новими областями (генеративний дизайн, ШІ-етика). Ключове значення має центральний конструкт TPACK~--- інтеграція всіх трьох вимірів, що забезпечує цілісне формування ІКК, а не фрагментарне опанування окремих технологій.

    \item \textit{Студійна педагогіка} як <<підписова педагогіка>> дизайнерських професій~\cite{Shulman2005}. Рефлексивна практика за Шоном~\cite{Schon1983}~--- рефлексія-в-дії та рефлексія-про-дію~--- є теоретичним підґрунтям рефлексивного компоненту ІКК. Оглядове дослідження~\cite{Musiienko2026ScopingReview} зафіксувало парадигмальний зсув від <<дизайнера-виробника>> до <<дизайнера-куратора>>, що потребує оновлення студійної педагогіки: критика (crit) має включати обґрунтування вибору інструментів та стратегії промптингу, ітеративний процес прискорюється завдяки генеративним моделям, формується нова компетенція курування (відбір, оцінка та редагування ШІ-результатів).

    \item \textit{Компетентнісний підхід} відповідно до стандарту В2 Дизайн та Європейської рамки кваліфікацій~\cite{Musiienko2025DesignEducation}. Аналіз прогалин стандарту В2 (п.~\ref{subsec:1.4}) визначив конкретні напрямки доповнення: інформаційно-аналітичний компонент (ЗК2)~--- критична оцінка ШІ-контенту; комунікативний (ЗК6, ФК5)~--- цифрові комунікаційні платформи; технологічний (ФК8, ФК12)~--- ШІ-інструменти та промпт-інженерія; рефлексивний (ЗК1, ЗК10)~--- етика ШІ у дизайні.
\end{enumerate}

\textbf{Дидактичні принципи.} Реалізація моделі забезпечується системою дидактичних принципів:

\begin{itemize}
    \item \textit{Принцип системності та наскрізності}: ШІ-інструменти інтегруються не як ізольований курс, а наскрізно через усі дисципліни навчального плану від 1-го до 4-го курсу. Цей принцип обґрунтований міжнародним бенчмаркінгом (п.~\ref{subsec:2.3}): найуспішніші програми (RCA, Aalto University) інтегрують ШІ-тематику наскрізно.

    \item \textit{Принцип доповнення, а не заміщення}: ШІ-інструменти доповнюють, але не замінюють традиційні навички (рисунок, живопис, макетування). Жодна з 7 проаналізованих міжнародних програм не замінює ручні навички ШІ (п.~\ref{subsec:2.3}). Цей принцип забезпечує збереження фундаментальної дизайнерської підготовки.

    \item \textit{Принцип прогресивності}: складність ШІ-інтеграції зростає поступово~--- від базової ШІ-грамотності (1 курс) через інструментальне використання (2--3 курси) до критичного та етичного осмислення (4 курс). Це відповідає моделі Aalto University, визначеній як найкраща практика (п.~\ref{subsec:2.3}).

    \item \textit{Принцип практикоорієнтованості}: формування ІКК відбувається через реальну проєктну діяльність у студійному форматі, а не через теоретичне вивчення технологій. Студійна педагогіка як <<навчання через діяльність>> (learning by doing) є найбільш ефективним середовищем для формування ІКК.

    \item \textit{Принцип рефлексивності}: кожна взаємодія з ШІ-інструментом супроводжується рефлексивним аналізом (етичним, естетичним, функціональним). Рефлексивний компонент ІКК, виявлений як найменш розвинений (п.~\ref{subsec:2.2}), потребує цілеспрямованого формування.

    \item \textit{Принцип відповідності ринку праці}: зміст ІКК-підготовки узгоджується з реальними вимогами роботодавців та актуальними інструментами ринку. Опитування роботодавців (п.~\ref{subsec:2.2}) підтвердило необхідність орієнтації на конкретні технології та навички.
\end{itemize}

\subsubsection{Змістовий блок}\label{sssec:3.1.3}

Змістовий блок визначає зміст формування ІКК через відображення 4 компонентів ІКК на навчальні дисципліни та програмні результати навчання. Цей блок є центральним у моделі, оскільки саме зміст навчання є безпосереднім об'єктом трансформації.

\textbf{Інформаційно-аналітичний компонент ІКК} передбачає формування здатності дизайнера до:
\begin{itemize}
    \item ефективного пошуку та систематизації візуальних референсів з використанням як традиційних баз даних (Behance, Dribbble, Pinterest), так і ШІ-засобів (пошук за зображенням, семантичний пошук);
    \item критичної оцінки достовірності, якості та етичності ШІ-генерованого контенту;
    \item аналізу даних для прийняття дизайнерських рішень (користувацькі дослідження, A/B тестування, аналітика);
    \item оцінки можливостей та обмежень різних ШІ-інструментів для конкретних дизайнерських завдань.
\end{itemize}

Цей компонент реалізується через інтеграцію аналітичних завдань у студійні курси (аналіз ШІ-референсів у Композиції, порівняльний аналіз результатів різних генеративних моделей у Графічному дизайні), а також через окрему дисципліну <<Дослідження в дизайні>> на 3 курсі.

\textbf{Комунікативний компонент ІКК} охоплює:
\begin{itemize}
    \item здатність до візуальної презентації дизайн-проєктів з використанням цифрових засобів та ШІ-інструментів (генерація мокапів, візуалізація концепцій);
    \item ефективну комунікацію в розподілених командах через цифрові платформи (Figma, Miro, Slack, Notion);
    \item формулювання технічних завдань та промптів для ШІ-інструментів як форму професійної комунікації <<людина--ШІ>>;
    \item здатність до академічної та професійної комунікації іноземною мовою, включаючи роботу з англомовною документацією ШІ-інструментів.
\end{itemize}

Комунікативний компонент реалізується через командні проєкти у студійних курсах, публічні критики (crit sessions) з обов'язковим використанням цифрових засобів презентації, а також через дисципліни <<Іноземна мова>> та <<Академічне письмо>>.

\textbf{Технологічний компонент ІКК} включає:
\begin{itemize}
    \item практичне володіння традиційними цифровими інструментами дизайну (Adobe Creative Suite, Figma, Blender, 3ds Max, AutoCAD та ін.);
    \item практичне володіння ШІ-інструментами дизайну: генеративні моделі (Midjourney, DALL-E, Stable Diffusion, Adobe Firefly), ШІ-функції у професійних інструментах (Adobe Sensei, Figma AI, Canva AI), генеративне відео та анімація (RunwayML, Pika);
    \item навички промпт-інженерії: формулювання ефективних текстових запитів, ланцюговий промптинг, негативні промпти, стилістичні модифікатори;
    \item інтеграція ШІ у робочі процеси дизайну: від ідеації через генерацію варіантів до фінальної доробки та виробництва.
\end{itemize}

Технологічний компонент реалізується через спеціалізовані дисципліни (<<Цифрова грамотність та основи AI>>, <<Комп'ютерна графіка>>, <<Генеративний дизайн та AI>>), а також через інтеграцію ШІ-інструментів у всі студійні та проєктні курси.

\textbf{Рефлексивний компонент ІКК} передбачає:
\begin{itemize}
    \item здатність до самоаналізу власного рівня ІКК та визначення напрямків саморозвитку;
    \item критичну оцінку етичних аспектів використання ШІ у дизайні (авторське право, bias, екологічний вплив, заміщення людської праці);
    \item рефлексію щодо ролі дизайнера у контексті розвитку ШІ (<<виробник>> vs. <<куратор>>);
    \item здатність обґрунтовувати вибір інструментів (традиційних чи ШІ-базованих) для конкретних дизайнерських завдань.
\end{itemize}

Рефлексивний компонент реалізується через критики (crit sessions) з обов'язковим обґрунтуванням інструментального вибору, рефлексивні щоденники у студійних курсах, а також через дисципліну <<Критичний дизайн та етика AI>> на 4 курсі. Аналіз висновків НАЗЯВО (п.~\ref{subsec:2.2}) засвідчив, що рефлексивний компонент має найменше представництво (11--12 згадувань), що підтверджує необхідність його цілеспрямованого формування.

\textbf{Відображення компонентів ІКК на навчальні дисципліни} наведено у таблиці~\ref{tab:ikk-course-mapping}.

\begin{longtable}{@{}p{0.20\textwidth}p{0.25\textwidth}p{0.22\textwidth}p{0.25\textwidth}@{}}
\caption{Відображення компонентів ІКК на типи навчальних дисциплін}\label{tab:ikk-course-mapping} \\
\hline
\textbf{Компонент ІКК} & \textbf{Профільні дисципліни} & \textbf{Фахові дисципліни} & \textbf{Загальноосвітні} \\
\hline
\endfirsthead
\hline
\textbf{Компонент ІКК} & \textbf{Профільні дисципліни} & \textbf{Фахові дисципліни} & \textbf{Загальноосвітні} \\
\hline
\endhead
Інформаційно-аналітичний & Дослідження в дизайні; Критичний дизайн & Історія мистецтва; Дизайн-мислення & Академічне письмо; Іноземна мова \\
\addlinespace
Комунікативний & UX/UI дизайн; Портфоліо та профпрактика & Графічний дизайн; Ілюстрація & Іноземна мова; Академічне письмо \\
\addlinespace
Технологічний & Генеративний дизайн та AI; Комп'ютерна графіка; Figma & 3D-моделювання; Моушн-дизайн & Цифрова грамотність та основи AI \\
\addlinespace
Рефлексивний & Критичний дизайн та етика AI; Дипломне проектування & Композиція; Живопис (crit) & Філософія; Дизайн та підприємництво \\
\hline
\end{longtable}

\subsubsection{Організаційно-технологічний блок}\label{sssec:3.1.4}

Організаційно-технологічний блок визначає інструменти, методи та форми організації навчального процесу, що забезпечують реалізацію змістового блоку. Цей блок конкретизує \textit{як} організовано формування ІКК.

\textbf{ШІ-інструменти для дизайн-освіти.} На основі аналізу оглядового дослідження~\cite{Musiienko2026ScopingReview} та міжнародного бенчмаркінгу (п.~\ref{subsec:2.3}) визначено базовий набір ШІ-інструментів, рекомендованих для інтеграції у навчальний процес:

\begin{longtable}{@{}p{0.22\textwidth}p{0.20\textwidth}p{0.20\textwidth}p{0.30\textwidth}@{}}
\caption{ШІ-інструменти для дизайн-освіти}\label{tab:ai-tools} \\
\hline
\textbf{Інструмент} & \textbf{Категорія} & \textbf{Дисципліни} & \textbf{Компоненти ІКК} \\
\hline
\endfirsthead
\hline
\textbf{Інструмент} & \textbf{Категорія} & \textbf{Дисципліни} & \textbf{Компоненти ІКК} \\
\hline
\endhead
Midjourney & Генерація зображень & Графічний дизайн, Ілюстрація, Композиція & Технол., Інф.-ан., Рефлекс. \\
\addlinespace
DALL-E (ChatGPT) & Генерація зображень, текст & Графічний дизайн, Дизайн бренду & Технол., Комунік. \\
\addlinespace
Stable Diffusion & Генерація зображень (відкритий код) & Генеративний дизайн та AI & Технол., Інф.-ан. \\
\addlinespace
Adobe Firefly & Інтегрований у Adobe Suite & Комп'ютерна графіка, Ілюстрація & Технол. \\
\addlinespace
Figma AI & UX/UI автоматизація & UX/UI дизайн, Figma та прототипування & Технол., Комунік. \\
\addlinespace
RunwayML & Генеративне відео & Моушн-дизайн, Відеографія & Технол., Рефлекс. \\
\addlinespace
Canva AI & Доступний генеративний дизайн & Типографіка, Інфографіка & Технол. \\
\addlinespace
ChatGPT/Claude & Текстовий ШІ-асистент & Усі дисципліни (дослідження, аналіз) & Інф.-ан., Комунік., Рефлекс. \\
\hline
\end{longtable}

Важливо підкреслити, що конкретні інструменти можуть змінюватися з розвитком технологій. Тому модель орієнтована не на опанування конкретного продукту, а на формування \textit{трансферних навичок} (промпт-інженерія, критична оцінка ШІ-результатів, інтеграція ШІ у робочий процес), що залишаються актуальними незалежно від інструменту.

\textbf{Методи навчання} адаптовано з урахуванням ШІ-інтеграції:

\begin{itemize}
    \item \textit{Студійний метод (модифікований)}: класичний цикл <<ідея~$\to$~прототип~$\to$~критика~$\to$~вдосконалення>> доповнюється етапом ШІ-генерації та порівняльного аналізу. Студент створює версію вручну, генерує ШІ-варіанти, порівнює та обґрунтовує фінальне рішення. Це формує одночасно технологічний (робота з ШІ), інформаційно-аналітичний (порівняльний аналіз) та рефлексивний (обґрунтування вибору) компоненти ІКК.

    \item \textit{Проєктне навчання з ШІ-компонентом}: кожен студійний проєкт включає обов'язковий етап використання щонайменше одного ШІ-інструменту. Студенти документують свій ШІ-робочий процес (промпти, ітерації, відбір результатів) як частину проєктної документації.

    \item \textit{Критика (crit) з ШІ-виміром}: публічне обговорення студентських робіт доповнюється аналізом стратегії використання ШІ. Студент має відповісти на запитання: <<Чому обрано саме цей інструмент?>>, <<Які промпти використано?>>, <<Як оцінено ШІ-результати?>>, <<Які етичні аспекти враховано?>>. Це цілеспрямовано формує рефлексивний компонент ІКК.

    \item \textit{Рефлексивний щоденник}: протягом навчання студенти ведуть цифровий щоденник, де фіксують свій досвід взаємодії з ШІ-інструментами, аналізують ефективність різних стратегій та рефлексують щодо етичних питань.

    \item \textit{Портфоліо-орієнтоване навчання}: формування цифрового портфоліо, що демонструє як традиційні роботи, так і ШІ-проєкти, є наскрізним завданням протягом усього навчання. Портфоліо виступає як інтегративний інструмент оцінювання всіх 4 компонентів ІКК.

    \item \textit{Командні проєкти з розподіленою роботою}: студенти працюють у командах, використовуючи цифрові платформи для колаборації (Figma, Miro, Notion), що формує комунікативний компонент ІКК.
\end{itemize}

\textbf{Форми організації навчання:}
\begin{itemize}
    \item студійні заняття (основна форма для фахових дисциплін);
    \item лекції з інтерактивними ШІ-демонстраціями;
    \item практичні заняття у комп'ютерних лабораторіях;
    \item самостійна проєктна робота з ШІ-інструментами;
    \item навчальна та виробнича практика з обов'язковим ШІ-компонентом;
    \item онлайн-консультації та критики через цифрові платформи.
\end{itemize}

\subsubsection{Діагностично-результативний блок}\label{sssec:3.1.5}

Діагностично-результативний блок визначає систему оцінювання сформованості ІКК та очікувані результати реалізації моделі. Цей блок забезпечує зворотний зв'язок для всіх попередніх блоків та є умовою ефективного функціонування моделі.

\textbf{Критерії оцінювання} відповідають 4 компонентам ІКК:

\begin{longtable}{@{}p{0.20\textwidth}p{0.30\textwidth}p{0.42\textwidth}@{}}
\caption{Критерії та показники оцінювання ІКК}\label{tab:ikk-criteria} \\
\hline
\textbf{Критерій} & \textbf{Компонент ІКК} & \textbf{Показники} \\
\hline
\endfirsthead
\hline
\textbf{Критерій} & \textbf{Компонент ІКК} & \textbf{Показники} \\
\hline
\endhead
Інформаційно-аналітичний & Пошук, оцінка та синтез інформації & Ефективність пошуку референсів; якість критичної оцінки ШІ-контенту; обґрунтованість аналітичних висновків \\
\addlinespace
Комунікативний & Візуальна та цифрова комунікація & Якість візуальних презентацій; ефективність командної взаємодії; якість промптів як комунікації з ШІ \\
\addlinespace
Технологічний & Володіння цифровими та ШІ-інструментами & Кількість опанованих інструментів; якість ШІ-генерованих результатів; ефективність інтеграції ШІ у робочий процес \\
\addlinespace
Рефлексивний & Самоаналіз та етична рефлексія & Глибина рефлексії у щоденнику; обґрунтованість інструментального вибору; врахування етичних аспектів \\
\hline
\end{longtable}

\textbf{Рівні сформованості ІКК:}

\begin{enumerate}
    \item \textit{Початковий (репродуктивний)}: здобувач знає основні ШІ-інструменти, але використовує їх фрагментарно, без критичного аналізу; промпти формулює на базовому рівні; рефлексія формальна; комунікація обмежена традиційними засобами.

    \item \textit{Базовий (адаптивний)}: здобувач систематично використовує 2--3 ШІ-інструменти у своїй роботі; формулює ефективні промпти; здатний до базової критичної оцінки ШІ-результатів; використовує цифрові засоби комунікації; веде рефлексивний щоденник.

    \item \textit{Достатній (конструктивний)}: здобувач вільно інтегрує ШІ-інструменти у робочий процес, обирає оптимальний інструмент для конкретного завдання; застосовує складні промпт-стратегії; критично оцінює ШІ-результати з естетичної, функціональної та етичної точок зору; ефективно комунікує через цифрові платформи.

    \item \textit{Високий (творчий)}: здобувач інноваційно застосовує ШІ-інструменти для створення оригінальних дизайн-рішень; розробляє власні промпт-стратегії; глибоко рефлексує щодо ролі ШІ у дизайні; є наставником для інших студентів; формує професійне цифрове портфоліо з ШІ-проєктами.
\end{enumerate}

\textbf{Інструменти оцінювання:}
\begin{itemize}
    \item \textit{Цифрове портфоліо} (основний інструмент): інтегративна оцінка всіх 4 компонентів ІКК через аналіз проєктних робіт, документації ШІ-процесів та рефлексивних записів. Портфоліо оцінюється за критеріями: технічна якість, креативність, обґрунтованість інструментального вибору, документація процесу, етична рефлексія.
    \item \textit{Рефлексивний щоденник}: оцінка рефлексивного компоненту через аналіз регулярних записів про досвід роботи з ШІ.
    \item \textit{Публічна критика (crit)}: оцінка комунікативного та рефлексивного компонентів через якість презентації та обґрунтування проєктних рішень.
    \item \textit{Проєктна документація}: оцінка технологічного та інформаційно-аналітичного компонентів через аналіз документації ШІ-процесів (промпти, ітерації, відбір).
    \item \textit{Самооцінка}: періодичне самооцінювання за 4 компонентами ІКК з використанням стандартизованого опитувальника.
\end{itemize}

\textbf{Очікувані результати} реалізації моделі:
\begin{itemize}
    \item підвищення рівня сформованості ІКК випускників за всіма 4 компонентами;
    \item подолання розриву між українськими та міжнародними програмами у частині ШІ-інтеграції;
    \item підвищення задоволеності випускників та роботодавців рівнем ІКТ-підготовки;
    \item формування професійних цифрових портфоліо з ШІ-проєктами.
\end{itemize}

\subsubsection{Візуалізація моделі}\label{sssec:3.1.6}

Структурно-функціональна модель формування ІКК майбутніх дизайнерів представлена на рис.~\ref{fig:structural-model}.

\begin{figure}[htbp]
\centering
\begin{tikzpicture}[
    node distance=1.2cm and 0.5cm,
    every node/.style={font=\small},
]

% ============ BLOCK 1: МОТИВАЦІЙНО-ЦІЛЬОВИЙ ============
\node[dia/header=15cm, fill=blue!8] (block1) {\textbf{МОТИВАЦІЙНО-ЦІЛЬОВИЙ БЛОК}};

\node[dia/lbox=4.5cm, below left=0.8cm and 3.5cm of block1] (demand1) {%
\textbf{Ринкова потреба:}\\
--- Рекомендація роботодавців\\
щодо інтеграції ШІ\\
--- Оцінка ІКТ: 3,5--4,0/5};

\node[dia/lbox=4.5cm, below=0.8cm of block1] (demand2) {%
\textbf{Незадоволеність здобувачів:}\\
--- 33,3\% (ГД) не задоволені ІКТ\\
--- 85,7\% (ДС) оцінка 1--3/5\\
--- Зниження ІКК в ОПП: $-$35\%};

\node[dia/lbox=4.5cm, below right=0.8cm and 3.5cm of block1] (demand3) {%
\textbf{Нормативні вимоги:}\\
--- Стандарт В2 (ФК8, ФК12)\\
--- DigComp 2.2: ШІ-грамотність\\
--- 0\% ШІ у 122 програмах};

\node[dia/rbox=10cm, below=3.5cm of block1, fill=blue!5] (goal) {%
\textbf{Мета:} формування ІКК майбутніх дизайнерів шляхом системної\\
трансформації навчальних програм з інтеграцією ШІ-інструментів};

% ============ BLOCK 2: ТЕОРЕТИКО-МЕТОДОЛОГІЧНИЙ ============
\node[dia/header=15cm, fill=green!8, below=0.8cm of goal] (block2) {\textbf{ТЕОРЕТИКО-МЕТОДОЛОГІЧНИЙ БЛОК}};

\node[dia/lbox=4.5cm, below left=0.7cm and 3.5cm of block2] (tpack) {%
\textbf{TPACK}\\
TK = ШІ + цифрові інструменти\\
PK = студійні методи\\
CK = предметна область дизайну};

\node[dia/lbox=4.5cm, below=0.7cm of block2] (studio) {%
\textbf{Студійна педагогіка}\\
Рефлексивна практика (Шон)\\
Виробник $\to$ Куратор\\
Критика з ШІ-виміром};

\node[dia/lbox=4.5cm, below right=0.7cm and 3.5cm of block2] (compet) {%
\textbf{Компетентнісний підхід}\\
Стандарт В2 Дизайн\\
EQF відповідність\\
4 компоненти ІКК};

% ============ BLOCK 3: ЗМІСТОВИЙ ============
\node[dia/header=15cm, fill=yellow!12, below=3.2cm of block2] (block3) {\textbf{ЗМІСТОВИЙ БЛОК: 4 КОМПОНЕНТИ ІКК}};

\node[dia/rbox=3.2cm, below left=0.7cm and 4.5cm of block3, fill=yellow!5] (ia) {%
\textbf{Інформаційно-}\\
\textbf{аналітичний}\\
Пошук, оцінка,\\
синтез інформації\\
(у т.\,ч. ШІ)};

\node[dia/rbox=3.2cm, below left=0.7cm and 0.8cm of block3, fill=yellow!5] (comm) {%
\textbf{Комунікативний}\\
Візуальна та цифрова\\
комунікація,\\
промпт-інженерія\\
як комунікація};

\node[dia/rbox=3.2cm, below right=0.7cm and 0.8cm of block3, fill=yellow!5] (tech) {%
\textbf{Технологічний}\\
Цифрові та ШІ-\\
інструменти,\\
генеративний\\
дизайн};

\node[dia/rbox=3.2cm, below right=0.7cm and 4.5cm of block3, fill=yellow!5] (refl) {%
\textbf{Рефлексивний}\\
Самоаналіз,\\
ШІ-етика,\\
обґрунтування\\
вибору};

% ============ BLOCK 4: ОРГАНІЗАЦІЙНО-ТЕХНОЛОГІЧНИЙ ============
\node[dia/header=15cm, fill=orange!10, below=3.8cm of block3] (block4) {\textbf{ОРГАНІЗАЦІЙНО-ТЕХНОЛОГІЧНИЙ БЛОК}};

\node[dia/lbox=4.5cm, below left=0.7cm and 3.5cm of block4] (tools) {%
\textbf{ШІ-інструменти:}\\
Midjourney, DALL-E,\\
Stable Diffusion,\\
Adobe Firefly, Figma AI,\\
RunwayML, ChatGPT};

\node[dia/lbox=4.5cm, below=0.7cm of block4] (methods) {%
\textbf{Методи:}\\
Модифікована студія,\\
Проєктне навчання з ШІ,\\
Критика з ШІ-виміром,\\
Рефлексивний щоденник};

\node[dia/lbox=4.5cm, below right=0.7cm and 3.5cm of block4] (forms) {%
\textbf{Форми:}\\
Студійні заняття,\\
Лабораторні роботи,\\
Командні проєкти,\\
Практика на підприємствах};

% ============ BLOCK 5: ДІАГНОСТИЧНО-РЕЗУЛЬТАТИВНИЙ ============
\node[dia/header=15cm, fill=red!8, below=3.2cm of block4] (block5) {\textbf{ДІАГНОСТИЧНО-РЕЗУЛЬТАТИВНИЙ БЛОК}};

\node[dia/lbox=4.5cm, below left=0.7cm and 3.5cm of block5] (criteria) {%
\textbf{Критерії:}\\
4 критерії за компонентами\\
ІКК (інф.-ан., комунік.,\\
технол., рефлекс.)};

\node[dia/lbox=4.5cm, below=0.7cm of block5] (instruments) {%
\textbf{Інструменти:}\\
Цифрове портфоліо,\\
Рефлексивний щоденник,\\
Публічна критика (crit),\\
Самооцінка};

\node[dia/lbox=4.5cm, below right=0.7cm and 3.5cm of block5] (levels) {%
\textbf{Рівні:}\\
Початковий (репродуктивний)\\
Базовий (адаптивний)\\
Достатній (конструктивний)\\
Високий (творчий)};

% ============ RESULT ============
\node[dia/header=15cm, fill=green!12, below=3.2cm of block5] (result) {%
\textbf{РЕЗУЛЬТАТ:} Випускник-дизайнер зі сформованою ІКК\\
(інформаційно-аналітичний + комунікативний + технологічний + рефлексивний)};

% ============ ARROWS ============
% Block 1 -> Goal
\draw[dia/arrow] (demand1.south) -- ++(0,-0.3) -| (goal.north west);
\draw[dia/arrow] (demand2.south) -- (goal.north);
\draw[dia/arrow] (demand3.south) -- ++(0,-0.3) -| (goal.north east);

% Goal -> Block 2
\draw[dia/arrow] (goal.south) -- (block2.north);

% Block 2 -> components
\draw[dia/arrow] (block2.south) -- (tpack.north);
\draw[dia/arrow] (block2.south) -- (studio.north);
\draw[dia/arrow] (block2.south) -- (compet.north);

% Components -> Block 3
\draw[dia/arrow] (tpack.south) -- ++(0,-0.3) -| ([xshift=-2cm]block3.north);
\draw[dia/arrow] (studio.south) -- ++(0,-0.4) -- (block3.north);
\draw[dia/arrow] (compet.south) -- ++(0,-0.3) -| ([xshift=2cm]block3.north);

% Block 3 -> IKK components
\draw[dia/arrow] (block3.south) -- (ia.north);
\draw[dia/arrow] (block3.south) -- (comm.north);
\draw[dia/arrow] (block3.south) -- (tech.north);
\draw[dia/arrow] (block3.south) -- (refl.north);

% IKK -> Block 4
\draw[dia/arrow] (ia.south) -- ++(0,-0.3) -| ([xshift=-3cm]block4.north);
\draw[dia/arrow] (comm.south) -- ++(0,-0.5) -| ([xshift=-1cm]block4.north);
\draw[dia/arrow] (tech.south) -- ++(0,-0.5) -| ([xshift=1cm]block4.north);
\draw[dia/arrow] (refl.south) -- ++(0,-0.3) -| ([xshift=3cm]block4.north);

% Block 4 -> sub-nodes
\draw[dia/arrow] (block4.south) -- (tools.north);
\draw[dia/arrow] (block4.south) -- (methods.north);
\draw[dia/arrow] (block4.south) -- (forms.north);

% Sub-nodes -> Block 5
\draw[dia/arrow] (tools.south) -- ++(0,-0.3) -| ([xshift=-2cm]block5.north);
\draw[dia/arrow] (methods.south) -- ++(0,-0.4) -- (block5.north);
\draw[dia/arrow] (forms.south) -- ++(0,-0.3) -| ([xshift=2cm]block5.north);

% Block 5 -> assessment nodes
\draw[dia/arrow] (block5.south) -- (criteria.north);
\draw[dia/arrow] (block5.south) -- (instruments.north);
\draw[dia/arrow] (block5.south) -- (levels.north);

% Assessment -> Result
\draw[dia/arrow] (criteria.south) -- ++(0,-0.3) -| ([xshift=-2cm]result.north);
\draw[dia/arrow] (instruments.south) -- ++(0,-0.4) -- (result.north);
\draw[dia/arrow] (levels.south) -- ++(0,-0.3) -| ([xshift=2cm]result.north);

% Feedback arrow (right side)
\draw[dia/darrow, dashed, gray] ([xshift=0.5cm]result.east) -- ++(1.5,0) |- ([xshift=0.5cm]block1.east)
    node[pos=0.5, right, text width=2cm, font=\footnotesize\itshape] {Зворотний\\зв'язок};

\end{tikzpicture}
\caption{Структурно-функціональна модель формування ІКК майбутніх дизайнерів}
\label{fig:structural-model}
\end{figure}

Модель функціонує як цілісна система зі зворотним зв'язком: результати діагностики (блок~5) впливають на коригування мети та завдань (блок~1), оновлення змісту (блок~3) та методів (блок~4). Зворотний зв'язок забезпечується через систематичний моніторинг сформованості ІКК засобами портфоліо-оцінювання та самооцінки.

\subsubsection{Характеристика зв'язків між блоками моделі}\label{sssec:3.1.7}

Функціонування моделі забезпечується двома типами зв'язків: вертикальними (ієрархічними) та горизонтальними (координаційними).

\textit{Вертикальні зв'язки} утворюють логічну послідовність: соціальне замовлення та мета (блок~1) $\to$ теоретичне обґрунтування та дидактичні принципи (блок~2) $\to$ конкретний зміст за 4 компонентами ІКК (блок~3) $\to$ інструменти, методи та форми реалізації (блок~4) $\to$ діагностика та корекція (блок~5). Ця послідовність забезпечує логічну зумовленість кожного наступного блоку попереднім.

\textit{Горизонтальні зв'язки} забезпечують координацію компонентів всередині кожного блоку. Зокрема:
\begin{itemize}
    \item у змістовому блоці: 4 компоненти ІКК взаємопов'язані та формуються паралельно (наприклад, промпт-інженерія є одночасно технологічною та комунікативною компетенцією);
    \item в організаційно-технологічному блоці: інструменти, методи та форми координуються для забезпечення інтегрального формування всіх компонентів ІКК;
    \item у діагностично-результативному блоці: критерії, інструменти та рівні утворюють єдину систему оцінювання.
\end{itemize}

\textit{Зворотний зв'язок} забезпечує саморегуляцію моделі: результати діагностики інформують про необхідність коригування на кожному рівні. Наприклад, якщо портфоліо-оцінювання виявляє недостатній рівень рефлексивного компоненту, це сигналізує про необхідність посилення рефлексивних завдань в організаційно-технологічному блоці та, можливо, перегляду змістового наповнення рефлексивного компоненту.

Запропонована структурно-функціональна модель відрізняється від існуючих моделей формування ІКТ-компетентностей кількома принциповими аспектами. По-перше, вона побудована на \textit{емпіричній основі}: мотиваційно-цільовий блок спирається не на загальні декларації, а на конкретні дані аналізу 122 програм, 132 силабусів, опитувань випускників, роботодавців та студентів. По-друге, модель інтегрує \textit{ШІ-вимір} як наскрізний елемент усіх блоків, а не як ізольований технологічний компонент. По-третє, рефлексивний компонент виокремлено як самостійний, з власними критеріями та інструментами оцінювання, що відповідає виявленій потребі у його цілеспрямованому формуванні~\cite{Musiienko2026ScopingReview}.

Конкретизацію змістового та організаційно-технологічного блоків моделі через рекомендації щодо трансформації навчальних дисциплін представлено у п.~\ref{subsec:3.2}, а педагогічні умови та стратегію впровадження~--- у п.~\ref{subsec:3.3}.
